\section{引言：大规模分布式训练的时代}

本节课聚焦于大规模分布式训练——这是当今所有神经网络在实践中训练的方式。10 年前，在一块 GPU 上训练模型是常态；而今天，在数千甚至数万块 GPU 上并行训练已成为新常态。课程以 \textbf{Llama 3 405B} 作为贯穿全课的案例研究，因为它是少数公开了详细训练基础设施信息的前沿模型之一。

\begin{figure}[H]
\centering
\includegraphics[width=0.95\textwidth]{slides-images/slide-001.jpg}
\caption{课程标题页：大规模分布式训练\protect\footnotemark}
\end{figure}
\footnotetext{Slides 第1页。}

\begin{importantbox}{为什么选择 Llama 3 作为案例}
OpenAI 的 GPT-4 论文明确声明``不再分享架构、模型规模、硬件、训练计算、数据集构建、训练方法等任何细节''。这已成为行业新常态。Llama 3 虽非最强模型，但其技术报告详细公开了训练集群和系统基础设施信息，为我们提供了一窥大规模 LLM 训练的珍贵窗口。
\end{importantbox}

\begin{figure}[H]
\centering
\includegraphics[width=0.95\textwidth]{slides-images/slide-002.jpg}
\caption{Llama 3 405B 模型概述与 GPT-4 的保密声明\protect\footnotemark}
\end{figure}
\footnotetext{Slides 第2页。}

课程分为两大部分：
\begin{enumerate}
    \item \textbf{Part 1}：GPU 硬件——理解我们用来训练模型的物理设备
    \item \textbf{Part 2}：分布式训练算法——如何在大量 GPU 上训练一个神经网络
\end{enumerate}

\subsection{本章小结}

大规模分布式训练已成为深度学习的新范式。理解硬件和并行化算法是构建高效训练系统的必要前提。Llama 3 的开放性使其成为学习这些技术的理想案例。

%% ========== GPU 硬件 ==========

